\section{The least index at minimal cyclotomic ramification}\label{sec:minimal-index}
This section proves Theorem~\ref{thm:minimal-index} and the exact characterization in Corollary~\ref{cor:exact-index-two}. The upper bound compares the graph's pointwise valuation constraint with the average valuation of primitive torsion. The lower bound is a finite twist of a Lubin--Tate group; its residual representations are identified by the Frobenius of a quadratic finite field.

\begin{proof}[The upper bound in Theorem~\ref{thm:minimal-index}]
Suppose $G[p^2]\simeq H[p^2]$ and $e=p-1$. Pass to the completed maximal unramified base. A determinant in the filtration class is $\chi_p^d\epsilon$, where $d\in\{0,1,2\}$ and $\epsilon\in\mu_{p-1}$. The wild cyclotomic image modulo $p^2$ has order $p$: otherwise $K$ would contain $\mu_{p^2}$, contrary to its absolute index. Equality of the determinant truncations forces the two dimensions $d$ to agree, because a nonzero difference of absolute value at most two cannot be divisible by $p\ge5$. Reduction is injective on $\mu_{p-1}$, so the finite determinant factors also agree.

For dimension zero, the special character is the Brauer lift of the residual semisimplification; duality treats dimension two. In dimension one, a potentially ordinary group has a unique connected--\'etale filtration, which descends. Thus each group either is whole potentially supersingular or is filtered by two potentially good height-one groups. Two filtered groups are compared residually, since $\chi_p\equiv1\pmod p$. A mixed comparison would require $e>p$ by Proposition~\ref{prop:mixed-obstruction} with $n=2$ and $v_p(e)=0$, which is impossible.

If both groups are potentially supersingular and their special characters differ, Proposition~\ref{fc-thm:valuation} gives
\[
 (p-1)(p^2-1)\ge(p-1)p^2.
\]
This contradiction proves recovery from $p^2$-torsion throughout the filtration class. The strict inequality comes from the graph bound on every point, not from a choice of Newton slope.
\end{proof}

\subsection{A pair with identical first-level torsion}
Put $k=\overline{\F}_p$, $R_0=W(k)$, $K_0=\operatorname{Frac}R_0$ and
\[
 \mathcal O=W(\F_{p^2}),\qquad q=p^2,
 \qquad K=K_0(\zeta_p).
\]
The cyclotomic polynomial $\Phi_p(1+X)$ is Eisenstein over $R_0$, so $K$ is complete, has residue field $k$, contains $\mu_p$ and has absolute index $p-1$.

Apply \citet[Lemma~1 and Theorem~1, pp.~381--382]{LT1965} over the unramified quadratic local field $\operatorname{Frac}\mathcal O$, with uniformizer $p$ and the series
\[
                         [p](T)=pT+T^q.
\]
This gives a one-dimensional formal $\mathcal O$-module, which we then base change to $R_0$. Its kernels of $[p^n]$ are finite flat of rank $q^n$, and their transition sequences define a height-two Barsotti--Tate group $G_0$. Its reduction is supersingular because $[p](T)$ reduces to $T^{p^2}$. The theorem is applied over a field with finite residue field before this base change.

The Tate module is free of rank one over $\mathcal O$ \citep[Theorem~2(b)--(d), pp.~383--384]{LT1965}. One can also check this directly: a primitive $p^n$-point has annihilator $p^n\mathcal O$, so the $q^n$ elements of its cyclic $\mathcal O$-module exhaust the $p^n$-torsion. Compatible primitive points give the assertion on inverse limits. All $\mathcal O$-endomorphisms are defined over $K_0$, so Galois acts through a character
\[
                       \psi:G_{K_0}\longrightarrow\mathcal O^*.
\]

Choose $t$ with $t^{q-1}=-p$. The Eisenstein polynomial $T^{q-1}+p$ gives a cyclic tame extension $E=K_0(t)$ of degree $q-1$, since $K_0$ contains $\mu_{q-1}$. Its nonzero roots are exactly the first-level torsion coordinates. We verify the subfield $K\subset E$ explicitly. If $\lambda=\zeta_p-1$, the cyclotomic equation gives
\[
 \lambda^{p-1}=-pu,\qquad
 u=1+\sum_{i=1}^{p-2}\frac{\binom p{i+1}}p\lambda^i
                      \in1+\lambda\mathcal O_K.
\]
Hensel's lemma supplies $v\in1+\lambda\mathcal O_K$ with $v^{p-1}=u$. Hence $(\lambda/v)^{p-1}=-p$. This root generates $K$ over $K_0$; all roots differ by elements of $\mu_{p-1}\subset K_0$. Since $t^{p+1}$ is another such root,
\[
                    K=K_0(t^{p+1}),\qquad [E:K]=p+1.
\]
For $a\in\mu_{q-1}$, the series $aT$ commutes with $pT+T^q$. Uniqueness in the Lubin--Tate construction gives $[a](T)=aT$. In the $\mathcal O/p$-basis $t$ of $G_0[p]$, the character $\overline\psi$ is therefore the reduction of $g(t)/t$. It follows that
\[
                 \overline\psi(G_K)=\mu_{p+1}\subset\F_{p^2}^*.
\]

Let $a(g)$ be the Teichm\"uller lift of $\overline\psi(g)$, and define
\[
                   \eta(g)=a(g)^{p-1}.
\]
The character $\eta$ has order $(p+1)/2\ge3$ and satisfies $\eta^p=\eta^{-1}$. Let $K'/K$ be the cyclic tame extension killing $\eta$. Twist the generic fibre of $G_0$ by the automorphisms $[\eta]$, obtaining $G_1$ with Tate character $\eta\psi$.

Here finite-level descent makes the construction explicit. For $M_n=\mathcal O/p^n\mathcal O$ with its $G_K$-action through $\eta\psi$, choose a finite Galois extension $D_n/K$ through which that action factors. The algebra
\[
 B_n=\operatorname{Map}(M_n,D_n)^{\operatorname{Gal}(D_n/K)},
 \qquad (gf)(x)=g\bigl(f(g^{-1}x)\bigr)
\]
is a product of fixed fields of point stabilizers and is finite \'etale over $K$. Addition, zero and inverse on $M_n$ give its Hopf operations. The equivariant maps $M_n\to M_{n+1}$, $x\mapsto px$, and $M_{n+1}\to M_n$, $x\mapsto x\bmod p^n$, give the closed embeddings and multiplication maps. Their kernel sequences are exact after separable base change and hence over $K$. Thus $\operatorname{Spec}B_n$ form a height-two Barsotti--Tate group. Over $K'$, the twist is trivial at every level, so $G_1$ has the same good supersingular model as $G_0$ after that extension.

Modulo $p$ one has $\overline\eta\,\overline\psi=\overline\psi^{p}$. The invertible $\F_p$-linear map
\[
                  \F_{p^2}\longrightarrow\F_{p^2},
                         \qquad x\longmapsto x^p
\]
intertwines the two Galois actions. It therefore gives $G_0[p]\simeq G_1[p]$ as finite group schemes over $K$.

The special character of $G_0$ is $2$, because it has good reduction over $K$. On the special isocrystal of $G_1$, the two embeddings of $\mathcal O$ give one-dimensional eigenspaces: Frobenius interchanges them, and their dimensions sum to two. The finite descent trace is consequently $\eta+\eta^p=\eta+\eta^{-1}$. A nontrivial value $a$ cannot satisfy $a+a^{-1}=2$, since that would imply $(a-1)^2=0$ in characteristic zero. The special characters differ. This proves the lower bound in Theorem~\ref{thm:minimal-index}, even within the class of whole potentially supersingular groups.

\begin{proof}[Proof of Corollary~\ref{cor:exact-index-two}]
For $p<e<p^3$, Theorem~\ref{thm:prime-index-upper} gives recovery from $p^3$-torsion, while Corollary~\ref{cor:sharp-mixed-prime} with $n=2$ gives unequal-character pairs with isomorphic $p^2$-torsion. Thus the least index is two.

Among positive multiples of $p-1$ with $e\le p$, the only possibility is $e=p-1$, where Theorem~\ref{thm:minimal-index} gives index one. Equality $e=p^3$ is excluded by $p\nmid e$. For $e>p^3$, the same mixed construction with $n=3$ gives failure of recovery from $p^3$-torsion, hence index at least three. These cases exhaust the complement of the displayed interval. The index is finite for every permitted $e$, since one can choose $n\ge3$ with $e<p^n$ and apply Theorem~\ref{thm:prime-index-upper}.
\end{proof}
